%% ============================================================
%%  SOLUCIÓN — Ejercicio 1.3: Teoremas y proposiciones
%%  Clase 2 — Después de diapositiva "Entornos de teoremas"
%%  Compilar: F5 (pdflatex)
%% ============================================================

\documentclass[12pt, a4paper]{article}
\usepackage[utf8]{inputenc}
\usepackage[T1]{fontenc}
\usepackage[spanish]{babel}
\usepackage{amsmath, amssymb, amsthm}

% --- ENTORNOS MATEMÁTICOS ---
\newtheorem{theorem}{Teorema}[section]
\newtheorem{proposition}[theorem]{Proposición}
\newtheorem{lemma}[theorem]{Lema}

\theoremstyle{definition}
\newtheorem{definition}{Definición}[section]
\newtheorem{assumption}{Supuesto}

\theoremstyle{remark}
\newtheorem{remark}{Observación}

\title{Marcos teóricos en econometría}
\author{Tu Nombre}
\date{\today}

\begin{document}
\maketitle

\section{Capital humano}

\begin{definition}
  El \textbf{estimador de Máxima Verosimilitud} (EMV) de un parámetro 
  $\theta \in \Theta$ es el valor $\hat{\theta}$ que maximiza la función 
  de verosimilitud:
  \[
    \hat{\theta} = \arg\max_{\theta \in \Theta} L(\theta; y_1, \ldots, y_n)
  \]
  donde $L(\theta; y_1, \ldots, y_n)$ es la verosimilitud de la muestra 
  observada bajo el parámetro $\theta$.
\end{definition}

\begin{assumption}
  \label{ass:gauss_markov}
  \textbf{Supuestos de Gauss-Markov}: Sea el modelo de regresión lineal 
  $Y = X\beta + \varepsilon$. Asumimos:
  \begin{enumerate}
    \item Linealidad en los parámetros.
    \item Muestra aleatoria.
    \item Variabilidad en $X$ (rango completo).
    \item Exogeneidad: $\mathbb{E}[\varepsilon | X] = 0$.
    \item Homocedasticidad: $\text{Var}(\varepsilon | X) = \sigma^2 I_N$.
    \item No autocorrelación: $\text{Cov}(\varepsilon_i, \varepsilon_j | X) = 0$ 
          para $i \neq j$.
  \end{enumerate}
\end{assumption}

\begin{proposition}
  \label{prop:ols_meli}
  Bajo los Supuestos de Gauss-Markov (Supuesto \ref{ass:gauss_markov}), 
  el estimador OLS $\hat{\beta} = (X'X)^{-1}X'Y$ es \textbf{MELI} 
  (Mejor Estimador Lineal Insesgado). Esto significa que entre todos los 
  estimadores que son lineales en $Y$ e insesgados, OLS tiene la varianza 
  más pequeña.
\end{proposition}

\begin{proof}
  El estimador OLS es:
  \[
    \hat{\beta} = (X'X)^{-1}X'Y = (X'X)^{-1}X'(X\beta + \varepsilon) 
              = \beta + (X'X)^{-1}X'\varepsilon
  \]
  
  Tomando esperanza:
  \[
    \mathbb{E}[\hat{\beta}] = \beta + (X'X)^{-1}X'\mathbb{E}[\varepsilon]
                            = \beta + 0 = \beta
  \]
  
  donde usamos el Supuesto de exogeneidad $\mathbb{E}[\varepsilon] = 0$.
  
  Luego, la matriz de varianza-covarianza es:
  \[
    \text{Var}(\hat{\beta}) = (X'X)^{-1}X'\text{Var}(\varepsilon)X(X'X)^{-1}
                            = \sigma^2 (X'X)^{-1}
  \]
  
  donde usamos homocedasticidad. Por el Teorema de Gauss-Markov, esta varianza 
  es la mínima entre todos los estimadores lineales insesgados.
\end{proof}

\begin{remark}
  La Proposición \ref{prop:ols_meli} es fundamental en econometría. Nos dice 
  que si los supuestos se cumplen, no hay mejor estimador que OLS (en la clase 
  de estimadores lineales e insesgados). Sin embargo, si alguno de los supuestos 
  falla —particularmente la exogeneidad— entonces OLS puede ser sesgado y 
  necesitamos otros métodos como VI (variables instrumentales).
\end{remark}

\end{document}
